\begin{figure}[htbp]
\centering
\setlength{\unitlength}{1mm}
\begin{picture}(110,82)
  \qbezier(10,2)(55,2)(100,2)
  \qbezier(100,2)(108,2)(108,10)
  \qbezier(108,10)(108,24)(108,38)
  \qbezier(108,38)(108,46)(100,46)
  \qbezier(100,46)(55,46)(10,46)
  \qbezier(10,46)(2,46)(2,38)
  \qbezier(2,38)(2,24)(2,10)
  \qbezier(2,10)(2,2)(10,2)
  \put(6,22){\makebox(0,0)[l]{$\Lambda=N[v]$}}
  \qbezier(55,8)(36,22)(17,36)
  \qbezier(55,8)(55,22)(55,36)
  \qbezier(55,8)(74,22)(93,36)
  \qbezier(17,36)(17,48.5)(17,61)
  \qbezier(17,36)(36,48.5)(55,61)
  \qbezier(55,36)(36,48.5)(17,61)
  \qbezier(55,36)(55,48.5)(55,61)
  \qbezier(93,36)(93,48.5)(93,61)
  \multiput(17,64)(0,2.5){5}{\line(0,1){1.2}}
  \multiput(55,64)(0,2.5){5}{\line(0,1){1.2}}
  \multiput(93,64)(0,2.5){5}{\line(0,1){1.2}}
  \put(55,8){\circle*{1.8}}
  \put(59,7){\makebox(0,0)[l]{$v$}}
  \put(17,36){\circle*{1.8}}
  \put(13,36){\makebox(0,0)[r]{$u_1$}}
  \put(55,36){\circle*{1.8}}
  \put(59,36){\makebox(0,0)[l]{$u_2$}}
  \put(93,36){\circle*{1.8}}
  \put(97,36){\makebox(0,0)[l]{$u_3$}}
  \put(17,61){\circle*{1.8}}
  \put(13,61){\makebox(0,0)[r]{$w_1$}}
  \put(55,61){\circle*{1.8}}
  \put(59,61){\makebox(0,0)[l]{$w_2$}}
  \put(93,61){\circle*{1.8}}
  \put(97,61){\makebox(0,0)[l]{$w_3$}}
  \put(55,80){\makebox(0,0){\small vertices outside $\Lambda$}}
\end{picture}
\caption{A root and its boundary. The outlined region is $\Lambda=N[v]$. The boundary vertices $w_1,w_2$ have the same set of neighbors in $N(v)$, namely $T_{w_1}=T_{w_2}=\{u_1,u_2\}$, but are recorded separately. The dotted lines indicate edges to vertices farther from the root. For each boundary vertex $w$, the flag $x_w$ records whether $w\in E$, and $b_w$ records whether all of $N[w]\setminus\Lambda$ lies in $E$.}
\label{fig:neighborhood}
\end{figure}
